\documentclass[12pt, a4paper]{article}
\usepackage[T1]{fontenc}
\usepackage[utf8]{inputenc}
\usepackage[polish]{babel}
\usepackage[margin=2.5cm]{geometry}
\usepackage{array}
\usepackage{tabularx}
\usepackage{setspace}
\usepackage{verbatim}
\usepackage[hidelinks]{hyperref}
\usepackage{tocloft}
\usepackage{float}
\renewcommand{\cftsecleader}{\cftdotfill{\cftdotsep}}

\onehalfspacing

\title{\vspace*{3cm} \textbf{Dokumentacja Funkcjonalna Aplikacji Java}}
\author{Autorzy: Mateusz Bombola \\ Kajetan Karwacki}
\date{\today}

\begin{document}

\maketitle
\thispagestyle{empty}

\newpage
\tableofcontents
\newpage

\section{Cel projektu}
Celem tej części projektu jest stworzenie okienkowej aplikacji w języku Java z wykorzystaniem biblioteki Swing. Aplikacja służy jako graficzny interfejs użytkownika do wizualizacji grafów planarnych. Głównym zadaniem programu jest wczytanie struktury grafu oraz jego współrzędnych (wygenerowanych wcześniej przez moduł obliczeniowy w języku C), a następnie interaktywne wyświetlenie ich na ekranie, dając użytkownikowi możliwość swobodnej manipulacji widokiem.

\section{Opis funkcji udostępnianych przez program}
Aplikacja została zaprojektowana z myślą o intuicyjnej obsłudze przy użyciu myszy i klawiatury. Główne funkcjonalności programu to:
\begin{itemize}
    \item \textbf{Wizualizacja grafu:} Rysowanie węzłów (reprezentowanych jako punkty) oraz krawędzi (linii łączących) na płótnie.
    \item \textbf{Interakcja z węzłami:} Możliwość ręcznej edycji i przesuwania pojedynczych węzłów po płótnie w celu poprawy czytelności układu.
    \item \textbf{Obsługa przestrzeni roboczej:} Przybliżanie i oddalanie widoku za pomocą rolki myszy lub przycisków w panelu.
    \item \textbf{Zarządzanie wyświetlaniem:} Pokaż/ukryj identyfikatory węzłów oraz wagi krawędzi.
    \item \textbf{Eksport wyników:} Zapis zmodyfikowanego ułożenia grafu do nowego pliku.
\end{itemize}

\section{Budowa Interfejsu Użytkownika}
Okno główne aplikacji podzielone jest na trzy sekcje:

\subsection{Pasek Menu}
Znajduje się na samej górze okna i zawiera opcje zarządzania plikami:
\begin{itemize}
    \item \textbf{Plik $\rightarrow$ Wczytaj strukturę grafu:} Pozwala wczytać początkowy plik tekstowy z definicją krawędzi.
    \item \textbf{Plik $\rightarrow$ Wczytaj współrzędne (.txt / .bin):} Wczytuje wygenerowane pozycje X i Y z pliku.
    \item \textbf{Plik $\rightarrow$ Zapisz wyniki (.txt / .bin):} Zapisuje obecne współrzędne węzłów do pliku.
    \item \textbf{Plik $\rightarrow$ Zapisz jako obraz:} Pozwala na zrzut obecnego widoku obszaru roboczego do pliku graficznego (np. PNG).
    \item \textbf{Plik $\rightarrow$ Zakończ:} Zamyka aplikację.
\end{itemize}

\subsection{Panel Narzędziowy}
Panel, w którym znajdują się przyciski szybkiego dostępu:
\begin{itemize}
    \item \textbf{Wybór algorytmu:} Rozwijana lista pozwalająca na wybór sposobu rozmieszczenia (Fruchterman-Reingold lub Tutte Embedding). Uwaga: Funkcja ta jest połączona z wywołaniem zewnętrznego programu C w tle.
    \item \textbf{Pokaż/Ukryj Etykiety:} Przełącznik sterujący widocznością numerów ID węzłów.
    \item \textbf{Pokaż/Ukryj Wagi:} Przełącznik sterujący widocznością wag na krawędziach.
    \item \textbf{Resetuj widok:} Przywraca domyślny poziom przybliżenia i środkuje graf na płótnie.
\end{itemize}

\subsection{Obszar Roboczy}
Centralna, największa część okna. Zaimplementowany komponent Swing, na którym za pomocą biblioteki \texttt{Graphics2D} renderowana jest ostateczna siatka grafu. Reaguje on na zdarzenia myszy (kliknięcia, przeciąganie, scrollowanie).

\section{Obsługiwane formaty danych}
Aplikacja wymaga do pełnego działania wczytania dwóch zestawów danych:

\subsection{Lista krawędzi (Struktura topologii)}
Plik tekstowy definiujący połączenia. Informuje program w Javie o tym, pomiędzy którymi węzłami należy narysować linie.
\\
\textbf{Format linii:} \texttt{<nazwa\_krawędzi> <wierzchołek\_A> <wierzchołek\_B> <waga>}
\\
\textbf{Przykład pliku wejściowego:}
\begin{verbatim}
AB 1 2 1.0
BC 2 3 1.5
CD 3 4 1.0
\end{verbatim}

\subsection{Współrzędne węzłów}
Aplikacja obsługuje oba formaty wygenerowane przez silnik obliczeniowy w języku C:
\begin{itemize}
    \item \textbf{Format tekstowy (.txt):} Parsowany linia po linii.
    \\ \textbf{Przykład pliku:}
\begin{verbatim}
1 0.0 0.0
2 100.5 50.0
3 50.0 100.0
\end{verbatim}
    \item \textbf{Format binarny (.bin):} Wczytywany za pomocą strumieni bajtowych. Program wczytuje początkową liczbę całkowitą (ilość węzłów), a następnie czyta kolejne 20-bajtowe bloki struktur (4 bajty ID typu Int, 8 bajtów osi X typu Double, 8 bajtów osi Y typu Double).
\end{itemize}
\section{Sytuacje wyjątkowe i kody błędów}
Aplikacja została zabezpieczona przed błędami użytkownika oraz uszkodzonymi danymi wejściowymi. Wszelkie anomalie nie powodują zawieszenia programu, lecz wyświetlają okno dialogowe z ostrzeżeniem.

\begin{table}[H]
\centering
\renewcommand{\arraystretch}{1.5}
\begin{tabular}{|p{4.5cm}|p{5cm}|p{5cm}|}
\hline
\textbf{Sytuacja wyjątkowa} & \textbf{Zachowanie programu} & \textbf{Komunikat dla użytkownika} \\
\hline
Wczytanie uszkodzonego pliku tekstowego & Przerwanie ładowania grafu, czyszczenie płótna. & \textit{"Błąd: Wybrany plik ma niepoprawny format lub zawiera litery zamiast liczb."} \\
\hline
Wczytanie pliku .bin niezgodnego ze strukturą & Zatrzymanie strumienia czytającego, odrzucenie danych. & \textit{"Błąd: Plik binarny ma niepoprawny format."} \\
\hline
Brak pliku topologii przed wczytaniem współrzędnych & Program ładuje punkty, ale nie rysuje krawędzi. & \textit{"Ostrzeżenie: Wczytano współrzędne, ale brakuje pliku krawędzi."} \\
\hline
Próba zapisu obrazu przy pustym płótnie & Blokada przycisku zapisu lub anulowanie akcji. & \textit{"Błąd: Brak grafu do zapisania. Obszar roboczy jest pusty"} \\
\hline
\end{tabular}
\end{table}

\end{document}
